 \section{解法の変更}\label{節：解法の変更}

  先に挙げた着目点から試した変形でも簡単にならない時は,一度最初からやり直してみるとよい.
  この際に今までと同じ考え方ではなく,なるべく別のアプローチができないか検討すると解法を見つけやすい.

  また,それ以外の場合でも次のどれかが起こったら,今の方法を一度止め,解法の再検討を行うと良い.
  \begin{enumerate}
    \item 対数を取っても,数列の項を1次の和として表せない.
    \item 定数を引いて付加項を消そうとしたが,条件を満たす定数が求まらない.
    \item 置換後に新しい種類の項が増え続ける.
  \end{enumerate}
  このようなときには$a_2$から$a_5$ぐらいまでの値を求めてみると法則性に気づき,帰納法で証明できる場合がある.
  

  それでも見つけられない時は,特殊関数型の置換や,まだ試していない式変形も候補に入れよう.
  一旦ご飯でも食べて改めて取り組もう.
  ご飯の時間でないのなら,今までの数学の知識を総動員して似たものを探そう.
  
  \begin{samepage}
  この手順を身につけると,個別のパターンを暗記するだけでなく,
  見慣れない漸化式にも最初の一手を選べるようになる.
  \par
  \end{samepage}
